%% kapa: exam=1 round=28 year=2017 period=1 form=A subject=real_estate subject_ko=부동산학원론 seq=34 number=114
%% source: Q-Net 감정평가사 기출문제 (cst00309) artlSeq=5199603 file=extracted/1008/5199603/2017년 제28회 감정평가사 1차시험 문제지/1교시 A형.pdf
%% conversion: kordoc 4.24.51 -> tools/gen1.py (auto); review=proofread-sonnet; flags=none
\begin{question}{114}
부동산 조세에 관한 설명으로 옳지 않은 것은?
\begin{choices}
  \item 재산세나 종합부동산세는 과세관청이 세액을 산정하여 납세의무자에게 교부하여 징수하는 세금인 반면, 상속세나 양도소득세는 납세의무자가 과세관청에 신고를 통해 납부하는 세금이다.
  \item 자본이득에 과세하는 양도소득세의 경우 소유자가 자산을 계속 보유함으로써 시장에서 자산거래가 위축되는 동결효과(lock-in effect)가 발생할 수 있다.
  \item 토지분 재산세의 과세대상 중 공장용지ㆍ전ㆍ답ㆍ과수원ㆍ목장용지와 같이 생산활동에 이용되는 토지는 별도 합산하여 과세한다.
  \item 취득세의 납세의무자는 사실상 취득자이다.
  \item 양도소득세의 양도가액은 원칙적으로 그 자산의 양도 당시의 양도자와 양수자 간에 실제로 거래한 가액에 따른다.
\end{choices}
\end{question}
\answer{3}
